\section{Zero in the Dolbeault spectrum}
\label{sec:spectral}

We first establish the analytic input needed for representations on Hilbert
spaces.  All metrics and differential forms on an orbifold are smooth in its
finite uniformizing charts.  In particular, the covering in the following
theorem need not be a manifold.

\begin{theorem}\label{thm:zero-spectrum}
Let $(M,\omega)$ be a connected compact K\"ahler orbifold of complex
dimension $n$, and let $U\to M$ be a connected normal orbifold covering with
infinite deck group $P$.  For some $p\in\{0,\ldots,n\}$, zero belongs to the
spectrum of the $L^2$ Dolbeault Laplacian on scalar $(p,0)$-forms on $U$,
with the lifted metric.
\end{theorem}

We use the nonnegative Laplacian
$\Delta''=\dbar\dbar^*+\dbar^*\dbar$.
For functions, the operator
$\tr_\omega\dd\dc$ has the opposite sign, up to a fixed positive
normalization.  Norms without a bundle subscript use the fixed bundle
metrics introduced below.

The invariant throughout the proof is the cohomology class of the
compact diagonal after push-down to the compact product. The estimates
below construct representatives with decreasing norm while controlling
the volume of their supports. We first build negative potentials and
a weighted inverse, then localize and smooth repeatedly. The final
pairing with the diagonal class makes these estimates incompatible
with a spectral gap in every degree.

\subsection{Complete operators and uniform local estimates}
\label{sp:operators}

Lifting a finite collection of nested uniformizing charts on $M$ gives
uniform coordinate radii, uniformly equivalent Euclidean metrics, bounds
for all metric derivatives, and partitions of unity with bounded overlap
and bounded derivatives.  A lifted chart has a subgroup of the original
finite uniformizing group.  The orders of all such groups are consequently
bounded.  Euclidean estimates applied to invariant lifts descend with
uniform constants, since the chart integrals differ from the orbifold
integrals by these bounded orders.  The same assertions hold on coverings
of $M\times M$ and for vector bundles with metrics pulled back from the
compact base.

We write $H^k$ for the real $L^2$ Sobolev spaces defined with these charts.
The local interior elliptic estimates, followed by summation over bounded
overlap, have the following consequence.  If $A$ is a scalar uniformly
elliptic second-order operator, or a system with scalar uniformly elliptic
principal symbol, whose coefficients have bounded derivatives of all
orders, then
\begin{equation}\label{sp:global-elliptic}
 \|v\|_{H^{k+2}}
 \le C_k\bigl(\|Av\|_{H^k}+\|v\|_2\bigr).
\end{equation}
Here the estimate applies whenever the right side is finite, initially in
the distributional sense.  One obtains it first on a smaller chart from
the corresponding larger chart and then sums the squared estimates.
The constants depend only on the ellipticity and the indicated coefficient
bounds; see the interior estimates in
\cite[Theorems~6.2 and~9.11]{GilbargTrudinger2001}, followed by
differentiation for higher Sobolev orders.
Sobolev embedding is uniform for the same reason.  In particular, an
$H^k$ function, with $k$ sufficiently large, tends to zero at infinity
together with any prescribed fixed number of derivatives: the local
Sobolev norms on fixed-radius charts outside a compact set are bounded by
the tail of its global Sobolev norm.

The lifted metric is complete.  Smooth compactly supported cutoffs
$\chi_R$, tending to one, can be chosen with
$|\dd\chi_R|\le C/R$ by smoothing truncated distance functions in the
uniform charts.  They show that the minimal and maximal closed extensions
of $\dbar$ agree.  They also justify the identities involving formal
adjoints: first approximate locally by smooth invariant sections and then
use
\[
 [\dbar,\chi_R]=\dbar\chi_R\wedge\ ,
 \qquad
 \|[\dbar,\chi_R]\|\le C/R.
\]
The same commutator bound holds for its adjoint and for the Chern
differential.  Thus compactly supported smooth sections form a core for
the associated quadratic forms.  Equivalently, the symmetric Dolbeault
Dirac operator $\sqrt2(\dbar+\dbar^*)$ is essentially self-adjoint.  For
completeness, its deficiency equation $D^*v=\pm\ii v$, tested against
$\chi_R^2v$, gives
$\|\chi_Rv\|_2^2\le C R^{-1}\|v\|_2^2$; local elliptic regularity
makes this test legitimate.  Letting $R$ tend to infinity eliminates both
deficiency spaces.  These arguments remain valid with each of the smooth
bundle multipliers used below, which have positive upper and lower bounds
for each fixed value of the parameters.  The K\"ahler identities
and the Bochner--Kodaira identity therefore extend from compact supports
to their complete-metric form domains; their local identities are those
of \cite[Chapter~VI, Section~6, and Chapter~VII, Theorems~1.1--1.2]{DemaillyCADG}.

Suppose, towards a contradiction to Theorem~\ref{thm:zero-spectrum}, that
all the scalar $(p,0)$ Laplacians have positive lower bounds.  Since there
are only finitely many degrees, choose a common lower bound $\delta>0$,
decreasing it whenever necessary by a fixed normalization constant.
The pointwise Lefschetz isomorphisms
\[
 L^{n-p}:\Lambda^{p,0}\longrightarrow\Lambda^{n,n-p},
 \qquad L=\omega\wedge\ ,
\]
commute with the scalar Laplacian and, after normalization, are isometries.
Thus the same type of positive lower bound holds on every $(n,q)$ degree.
In degree zero it gives
\begin{equation}\label{sp:poincare}
 \|v\|_2^2\le C\|\dd v\|_2^2,
 \qquad v\in H^1(U).
\end{equation}
The group $P$ is infinite, so $U$ is noncompact and $n>0$.

\subsection{Negative potentials with controlled mass}
\label{sp:potentials}

We need negative potentials with increasing depth on each fixed compact
set and controlled global \(L^2\) mass. The depth will make the later
weights large near that set, while the mass bound limits the volume of
regions where the potential is very negative and hence the growth of
supports during localization.

\begin{lemma}\label{sp:potential-lemma}
Under \eqref{sp:poincare}, there are smooth functions $u_N$ on $U$, for
all sufficiently large integers $N$, such that
\begin{equation}\label{sp:potential-bounds}
 u_N\le0,\qquad \dd\dc u_N\ge-\omega,
 \qquad \|u_N\|_2\le C N^n.
\end{equation}
For every compact set $K\subset U$, there are $c_K>0$ and $N_K$ with
\begin{equation}\label{sp:potential-depth}
 \sup_Ku_N\le-c_KN\qquad(N\ge N_K).
\end{equation}
For each fixed $N$, the function $u_N$ belongs to every $H^k$ and has
bounded derivatives of all orders.
\end{lemma}

\begin{proof}
Fix a nonnegative $f\in C_c^\infty(U)$ which is strictly positive on a
chart ball, and put
\[
 F_N=\log(1+N^nf),\qquad \lambda=N^{-2n}.
\]
We solve
\begin{equation}\label{sp:ma-equation}
 \omega_u^n=e^{F_N+\lambda u}\omega^n,
 \qquad \omega_u=\omega+\dd\dc u>0.
\end{equation}
We give the complete continuity argument, including the global estimates
which distinguish this problem from a bounded-solution existence result.

Fix an integer $k>n+6$ and consider
\[
 \mathcal B=\{u\in H^{k+2}(U,\R):
                 \omega_u\ge c\omega\text{ for some }c>0\}.
\]
This is open in $H^{k+2}$.  The map
$u\mapsto\log(\omega_u^n/\omega^n)-\lambda u$ is smooth from
$\mathcal B$ to $H^k$, by Sobolev multiplication and composition in the
uniform charts.  Its linearization at $u$ is
\begin{equation}\label{sp:ma-linearization}
 v\longmapsto\tr_{\omega_u}\dd\dc v-\lambda v.
\end{equation}
At a smooth solution with bounded derivatives this is an isomorphism
$H^{k+2}\to H^k$.  Indeed its negative is coercive on $H^1$ when
integrated with $\omega_u^n$: its quadratic form is a positive multiple
of the gradient norm plus $\lambda\|v\|_2^2$.  Lax--Milgram gives an
$H^1$ inverse, and \eqref{sp:global-elliptic} gives the asserted higher
regularity and bounded inverse.  The metric $\omega_u$ is complete and
uniformly equivalent to $\omega$ at this individual solution.  A solution
in $\mathcal B$ is smooth with bounded derivatives by local elliptic
bootstrapping, since initially $k>n+6$.  The implicit function theorem
therefore gives openness for
\begin{equation}\label{sp:ma-path}
 \log(\omega_u^n/\omega^n)-\lambda u=sF_N,
 \qquad 0\le s\le1.
\end{equation}
At $s=0$ the solution is $u=0$.

We next obtain estimates uniform in $s$, with $N$ fixed.  Every solution
in $\mathcal B$ tends to zero at infinity with its first several
derivatives.  At an interior maximum of $u$, the determinant ratio is at
most one; at an interior minimum it is at least one.  The same bounds
hold when an extremum is approached at infinity, where $u$ tends to zero.
Consequently
\begin{equation}\label{sp:ma-czero}
 -\lambda^{-1}\|F_N\|_\infty\le u\le0.
\end{equation}
In particular the determinant ratio in \eqref{sp:ma-path} has positive
upper and lower bounds depending only on $N$.

Here are the second-order and higher local estimates in the order in
which they are needed.  Write $\Delta_u=\tr_{\omega_u}\dd\dc$.
Since $\tr_\omega\dd\dc u\ge-n$,
\[
 \Delta_\omega(sF_N+\lambda u)\ge-C_N-\lambda n.
\]
The differential trace estimate
\cite[Lemma~2.2]{CampanaGuenanciaPaun2013} is
\[
 \Delta_u\log\tr_\omega\omega_u
 \ge\frac{\Delta_\omega(sF_N+\lambda u)}
            {\tr_\omega\omega_u}
       -B\tr_{\omega_u}\omega.
\]
It gives
\begin{equation}\label{sp:yau-trace}
 \Delta_u\log\tr_\omega\omega_u
 \ge-B\tr_{\omega_u}\omega-C_N.
\end{equation}
In this estimate $B$ depends on a lower bound for the background
holomorphic bisectional curvature.  The possible negative term involving
$\Delta_\omega(sF_N+\lambda u)$ is divided by
$\tr_\omega\omega_u$, which is bounded below by the determinant lower
bound.
Apply \eqref{sp:yau-trace} to
$Q=\log\tr_\omega\omega_u-Au$, where $A>B+1$ is fixed.  Since
$\Delta_u u=n-\tr_{\omega_u}\omega$, at a maximum of $Q$ one obtains
\[
 (A-B)\tr_{\omega_u}\omega\le An+C_N.
\]
If the supremum is not attained, it is bounded by the value $\log n$ at
infinity.  At an attained maximum, the determinant upper bound and the
inverse-trace bound control $\tr_\omega\omega_u$ there.  The oscillation
bound in \eqref{sp:ma-czero} then controls it everywhere.  It follows that
\begin{equation}\label{sp:ma-equivalence}
 c_N\omega\le\omega_u\le C_N\omega
\end{equation}
uniformly along the path.

This also bounds the background real Laplacian of $u$.  Local Poisson
$W^{2,p}$ estimates, for $p$ larger than the real dimension, give a
uniform $C^{1,\alpha}$ bound on smaller charts.  The equation now has a
uniformly elliptic complex Hessian and a controlled H\"older right side.
Fix $0<\beta<\alpha$.  The local complex Monge--Amp\`ere
$C^{2,\beta}$ estimate
\cite[arXiv:1111.0902v1, Theorem~1.1]{YuWang2012} applies after adding
a local potential for $\omega$: the resulting function is strictly
plurisubharmonic with local sup norm controlled by \eqref{sp:ma-czero},
its Euclidean Laplacian is bounded by \eqref{sp:ma-equivalence}, and
the positive density is bounded below with its $n$-th root uniformly
$C^\alpha$.  Differentiation of the equation and local Schauder
estimates then give bounds for every derivative.  These are local bounds
uniform over all lifted charts and
all $s$, with constants depending on $N$ and the derivative order; no
global higher Sobolev bound has been assumed in obtaining them.

There is a global estimate with a better dependence on $N$.  Integration
by parts gives
\begin{align}
 \int_U(-u)(\omega_u^n-\omega^n)
 &=\sum_{a=0}^{n-1}\int_U
   \ii\partial u\wedge\dbar u\wedge
   \omega_u^a\wedge\omega^{n-1-a},\label{sp:ma-energy}\\
 c\|\dd u\|_2^2
 &\le\int_U(-u)(\omega_u^n-\omega^n)
 \le N^n\|f\|_2\|u\|_2.\label{sp:ma-energy-bound}
\end{align}
The first mixed term is the fixed background metric, so the lower
constant is independent of $N$.  For the upper bound use $u\le0$ and
\[
 e^{sF_N+\lambda u}-1\le (1+N^nf)^s-1\le N^nf.
\]
The integrand in \eqref{sp:ma-energy} is absolutely integrable: away from
$\supp f$, its absolute value is at most $\lambda u^2$, and on
$\supp f$ it is integrable.  To justify the integration by parts, insert
the complete-metric cutoffs.  The mixed metrics are bounded by
\eqref{sp:ma-equivalence}, and the cutoff errors tend to zero by
$u,\dd u\in L^2$.  Combining \eqref{sp:ma-energy-bound} with
\eqref{sp:poincare} gives
\begin{equation}\label{sp:ma-ltwo}
 \|u\|_2\le C N^n
\end{equation}
uniformly in $s$ and $N$.

It remains to establish closedness in the global Banach space.  In local
coordinates, \eqref{sp:ma-path} can be written
\[
 a_u^{i\bar j}u_{i\bar j}-\lambda u=sF_N,
 \qquad
 a_u^{i\bar j}=\int_0^1(g_{a\bar b}+t u_{a\bar b})^{i\bar j}\,dt.
\]
The superscript $i\bar j$ denotes the corresponding entry of the inverse
Hermitian matrix.
The coefficients have uniform ellipticity and all required derivative
bounds by the preceding local estimates.  Thus
\eqref{sp:global-elliptic} and \eqref{sp:ma-ltwo} give global bounds in
every fixed Sobolev order.  If $s_\nu\to s$, a subsequence of the
corresponding solutions converges smoothly on compact sets.  The limit
solves \eqref{sp:ma-path}, satisfies the uniform positive lower bound in
\eqref{sp:ma-equivalence}, and belongs to $H^{k+2}$ by lower
semicontinuity of the global norms.  It is therefore still in
$\mathcal B$.  This proves closedness and hence solvability at $s=1$.
Denote that solution by $u_N$.  The first three assertions of
\eqref{sp:potential-bounds} and the stated regularity have been proved.

Suppose \eqref{sp:potential-depth} fails for a fixed nonempty compact set
$K$.  Along a subsequence,
$\sup_K(u_N/N)\to0$.  Local plurisubharmonic compactness, after adding
a local potential for $\omega/N$, gives an $L^1_{\mathrm{loc}}$ limit
$v$ which is plurisubharmonic and nonpositive.  The alternative of
uniform collapse to $-\infty$ is excluded by the supremum on $K$.
Hartogs' lemma implies that $v$ attains zero on $K$.  Since $U$ is
connected, the maximum principle gives $v=0$ everywhere.

The trace--determinant inequality applied to \eqref{sp:ma-equation}
gives
\begin{equation}\label{sp:ma-trace-limit}
 \tr_\omega\dd\dc(u_N/N)
 \ge n f^{1/n}e^{\lambda u_N/n}-n/N.
\end{equation}
Pass to a further almost-everywhere convergent subsequence of $u_N/N$.
As $n\ge1$,
$\lambda u_N=N^{1-2n}(u_N/N)\to0$ almost everywhere.  The right side of
\eqref{sp:ma-trace-limit} converges in $L^1$ on a ball where $f>0$ to
$nf^{1/n}>0$, by dominated convergence and $u_N\le0$.  The left side
converges to zero as a distribution because $u_N/N\to0$ in local
$L^1$.  Testing against a nonzero nonnegative function supported in that
ball is a contradiction.  This proves \eqref{sp:potential-depth}.
\end{proof}

Choose once and for all a compact region $L\subset U$, of positive
volume, whose translates have interiors covering $U$.  It may be
enlarged to be connected.  Properness and cocompactness of the deck action
give bounded overlap and a bounded number of neighbors within any fixed
distance.  For $a,g\in P$ put
\begin{equation}\label{sp:cell-depth}
 A_a(g)=-\sup_{x\in gL}u_N(a^{-1}x).
\end{equation}

\begin{lemma}\label{sp:cell-lemma}
There are $C,\epsilon>0$, independent of $N,a,g$, such that
\begin{equation}\label{sp:cell-count}
 \#\{g\in P:A_a(g)\ge1\}\le C N^{2n},
\end{equation}
and
\begin{equation}\label{sp:cell-exponential}
 \int_{gL}\exp\!\left(
       \frac{-\epsilon u_N(a^{-1}x)}{A_a(g)+1}\right)\omega^n\le C.
\end{equation}
For each $R<\infty$ there is $C_R$ such that
\begin{equation}\label{sp:cell-neighbors}
 C_R^{-1}(A_a(g)+1)\le A_a(h)+1\le C_R(A_a(g)+1)
\end{equation}
whenever the distance between $gL$ and $hL$ is at most $R$.
\end{lemma}

\begin{proof}
Every cell counted in \eqref{sp:cell-count} contributes at least
$\vol(L)$ to the integral of $|u_N(a^{-1}x)|^2$.  Bounded overlap and
\eqref{sp:potential-bounds} prove the count.

Translate $gL$ to $L$ and normalize the potential by $A_a(g)+1$.
The resulting functions are nonpositive, have $\dd\dc$ bounded below
by $-\omega$, and have supremum in $[-1,0]$ on $L$.  Plurisubharmonic
compactness on a connected neighborhood of any fixed enlargement of
$L$ makes this a relatively compact family in local $L^1$.  In
particular its supremum cannot tend to $-\infty$ on another fixed
cell.  There are only finitely many possible relative neighbors; this
gives the upper comparison in \eqref{sp:cell-neighbors}, and exchanging
the two cells gives the lower comparison.

On finitely many smaller charts covering $L$, add fixed local potentials
to make the normalized functions plurisubharmonic.  Each limit has a
positive local exponential integrability exponent.  Compactness and the
semicontinuity and integrability statement of
\cite[Main Theorem~0.2]{DemaillyKollar2001} give a common sufficiently
small exponent and a uniform integral bound.  The added local potentials
are bounded, so they can be removed at the cost of a constant.  Summing
the chart estimates and translating back gives
\eqref{sp:cell-exponential}.
\end{proof}

\subsection{A weighted inverse on the diagonal covering}
\label{sp:weighted}

The potentials now have both a depth estimate and a uniform exponential
integrability bound on cells. The next step is an inverse estimate
that makes a closed form inexpensive to solve where the potential is
deep. Its dependence on the weight, not merely existence of an inverse,
will be needed when the solution is localized.

Set
\[
 D=(U\times U)/P,\qquad \Omega_S=\omega_x+S\omega_y,
 \qquad E=p_y^*K_M,
\]
where $P$ acts diagonally and $S\ge1$.  The maps $p_x,p_y$ here are to
$M$.  We give $E$ the metric from the original, unscaled metric
$\omega$.  The space $D$ is an orbifold covering of $M\times M$;
write $m=2n$ for its complex dimension.  The subgroup defining this
covering is the inverse image of the diagonal in $P\times P$.
In particular the diagonal map $M\to M\times M$ lifts to $D$.

Unweighted norms for $\Omega_S$ are denoted by $\|\cdot\|_S$; an
additional bundle metric multiplier $h$ is indicated by
$\|\cdot\|_{S,h}$.  On the $(m,*)$ complex with coefficients $E^*$,
the unweighted Dolbeault Laplacian has a positive lower bound $\delta$
independent of $S$.  To verify this on a compactly supported section of
$D$, lift it to $U\times U$ and integrate first in $x$, with $y$ in a
fundamental region.  Product splitting expresses the energy as its
nonnegative $x$ and $y$ parts.  The $x$ part uses scalar $(n,q_x)$-forms,
so the previously obtained gaps apply.  The same argument on
$E^*$-valued $(p,0)$-forms uses the original scalar $(p_x,0)$ gaps and
gives
\begin{equation}\label{sp:x-gap}
 \|\dbar\xi\|_S^2\ge\delta\|\xi\|_S^2
 \quad\text{for $\xi$ of type $(p,0)$.}
\end{equation}
These calculations are local product calculations and hence descend to
the diagonal quotient.  Bounded finite chart stabilizers do not change
them.

\begin{lemma}\label{sp:weighted-inverse}
Suppose a smooth positive multiplier $h$ on $E^*$ satisfies
\begin{equation}\label{sp:weight-hypotheses}
 1\le h\le e^b,
 \qquad \ii\Theta_{E^*,h}\ge-C_0t\Omega_S,
 \qquad b,t\ge0.
\end{equation}
Then the complete $\dbar$ complex of $E$-valued $(0,*)$-forms with
multiplier $k=h^{-1}$ is exact with closed ranges.  Every closed form
$F$ in this complex has a solution $\dbar v=F$ satisfying
\begin{equation}\label{sp:weighted-solution}
 \|v\|_{S,k}\le
 C\bigl(e^{b/4}+\sqrt t\,e^{b/2}\bigr)\|F\|_{S,k}.
\end{equation}
The constant depends on $M$, $C_0$ and $\delta$, but not on $S,b,t$.
\end{lemma}

\begin{proof}
First work on the $(m,*)$ complex with coefficients $E^*$.  Its
unweighted gap gives exactness and closed ranges: the bounded Green
operator yields the decomposition into the images of $\dbar$ and
$\dbar^*$.  Weighted and unweighted differential graph norms are
equivalent for fixed $b$.  Thus the domains, kernels and ranges of the
closed differential $d=\dbar$ are the same, and weighted exactness and
closed ranges follow.  This argument does not assert equality of the
weighted and unweighted adjoint domains.

Let $D'_h$ be the $(1,0)$ part of the Chern connection and let $*_S$ be
the complex-linear Hodge star on form indices.  On $(m,q)$-forms it
takes values in $(m-q,0)$, and
\[
 (D'_h)^*=\pm *_S\dbar *_S.
\]
This formula acts trivially on the coefficient index.  The derivative of
the bundle metric in the adjoint cancels the corresponding Chern
connection term.  The Bochner--Kodaira identity and
\eqref{sp:weight-hypotheses} imply
\begin{equation}\label{sp:bochner-weighted}
 \mathcal E_{S,h}(\eta)+C_1t\|\eta\|_{S,h}^2
 \ge\|(D'_h)^*\eta\|_{S,h}^2
 \ge\delta\|\eta\|_S^2.
\end{equation}
Here
$\mathcal E_{S,h}(\eta)=
\|\dbar\eta\|_{S,h}^2+\|\dbar_h^*\eta\|_{S,h}^2$
is the weighted Dolbeault energy.
Indeed $D'_h\eta=0$ in top holomorphic degree, and the curvature
commutator there is a sum of $q$ curvature eigenvalues.  For the last
inequality drop $h\ge1$ from the differential norm and apply
\eqref{sp:x-gap} to $*_S\eta$.  The star is an isometry.  Compact
support approximation from Section~\ref{sp:operators} proves the same
inequality on the form domain.

Fix a degree $q$ and put $A=dd_h^*$ on the closed subspace
$\im d$.  For a nonzero vector $\eta$ in a spectral band of $A$ up to
$\mu>0$, its minimal weighted lift has squared norm
$\langle A^{-1}\eta,\eta\rangle_{S,h}$, hence at least
$\mu^{-1}\|\eta\|_{S,h}^2$.  The minimal unweighted lift is also an
admissible weighted lift, with norm at most
$e^{b/2}\delta^{-1/2}\|\eta\|_S$.  Consequently
\[
 \|\eta\|_S^2
 \ge\delta e^{-b}\mu^{-1}\|\eta\|_{S,h}^2.
\]
On this spectral band $d\eta=0$ and
$\mathcal E_{S,h}(\eta)\le\mu\|\eta\|_{S,h}^2$.
Substitution in \eqref{sp:bochner-weighted} gives
\begin{equation}\label{sp:spectral-quadratic}
 \mu(\mu+C_1t)\ge\delta^2e^{-b}.
\end{equation}
This holds for every nonzero lower spectral band.  Solving the quadratic
inequality shows that the minimal inverses of $d$ have norm at most
$C(e^{b/4}+\sqrt t\,e^{b/2})$ in every degree.

Finally use the conjugate-linear bundle-metric star for Serre duality,
which is distinct from the linear star used above.  It identifies the
$(m,q)$ complex for $E^*$ with multiplier $h$ with the reversed
$(0,m-q)$ complex for $E$ with multiplier $h^{-1}$, interchanging the
differential and its adjoint.  It is an isometry for these norms.  The
complete-domain identities therefore transfer exactness, closed ranges,
and the inverse bounds to the latter complex.  This proves
\eqref{sp:weighted-solution}.
\end{proof}

\subsection{A smoothing operator with fixed propagation}
\label{sp:smoothing}

For the rest of the proof all distances, cell neighborhoods, and support
volumes refer to the standard metric $\Omega_1$.  This convention is
independent of the parameter $S$ in the preceding estimates.
Write $\pi:D\to M\times M$ for the covering and
$E_0=p_y^*K_M$ for the bundle downstairs, so that $E=\pi^*E_0$.

\begin{lemma}\label{sp:smoothing-lemma}
There is a smoothing operator $K$ on the complete Dolbeault complex of
$E$-valued $(0,*)$-forms on $D$ with the following properties:
\begin{enumerate}[label=\textup{(\roman*)}]
\item $K$ commutes with $\dbar$ and preserves separately the
antiholomorphic degrees in $x$ and $y$.
\item Its propagation is at most a fixed $R$, and its kernel and all
kernel derivatives are uniformly bounded in the lifted charts.
\item For compactly supported currents, $K$ commutes with trace to
$M\times M$.  On closed such currents it preserves the cohomology class
of their trace.
\end{enumerate}
If $B$ is an $E$-valued $(0,n)$ current with cellwise $L^1$ coefficients,
then
\begin{equation}\label{sp:smoothing-norm}
 \|KB\|_S\le C S^{n/2}
 \left(\sum_C\left(\int_C|B|_S\,\dd V_1\right)^2\right)^{1/2},
\end{equation}
where $C$ runs through the product cells modulo the diagonal action.
\end{lemma}

\begin{proof}
Choose an even Schwartz function $f_0$ whose Fourier transform is smooth
and compactly supported, with $f_0(0)=1$, and set
$K=f_0(\sqrt{\Delta''_1})$, using the standard unweighted product
metric and the fixed metric on $E$.  The product splitting of
$\Delta''_1$ proves the degree assertions.  Commutation with $\dbar$
follows first on its domain from the Dolbeault identities and then on
currents by duality.

The wave equation has finite propagation with a speed bounded by the
principal symbol.  In the present setting this can be seen from the
symmetric first-order operator $D_1=\sqrt2(\dbar+\dbar^*)$:
the energy identity for a moving cutoff contains only its commutator,
whose norm is bounded by a fixed multiple of the cutoff gradient.
It gives the domain-of-dependence estimate for $e^{\ii tD_1}$.
Since $f_0$ is even, the Fourier expression for
$f_0(\sqrt{\Delta''_1})$ uses only these finite-time wave operators.
Thus $K$ has fixed propagation.  The coefficient curvature contributes
only lower-order terms and does not alter this argument.  The identical
energy calculation in invariant charts proves it on the orbifold.

To see smoothing with uniform kernel bounds, use
\eqref{sp:global-elliptic} and local Sobolev embedding to bound point
evaluation, and derivative evaluation, after $(1+\Delta''_1)^{-a}$
for $a$ sufficiently large.  The factorization
\[
 K=(1+\Delta''_1)^{-a}
 \bigl((1+\Delta''_1)^{2a}K\bigr)
 (1+\Delta''_1)^{-a}
\]
has bounded middle factor by the Schwartz decay of $f_0$.  Evaluating
the outside factors and their adjoints bounds the kernel at any two
points, uniformly in the charts.  Increasing $a$ gives every derivative
bound.  This argument also defines $K$ on distributions.

For a compactly supported $E$-valued current $A$, define its orbifold
trace by
\[
 \langle\operatorname{Tr}A,\beta\rangle
 =\langle A,\pi^*\beta\rangle
\]
for smooth $E_0^*$-valued test forms $\beta$ of complementary bidegree.
On a covering chart $[V/H]\to[V/G]$, this is the sum over cosets
$G/H$, with the natural action on form and coefficient indices, summed
also over the local sheets met by the support.  The $|G:H|$ terms
convert the downstairs integration factor $1/|G|$ into the upstairs
factor $1/|H|$.  This description applies even when $H$ is not normal;
it includes the stabilizer multiplicities at singular strata.

For a compactly supported current, the wave-evolved support stays in a
fixed compact neighborhood for bounded time, by completeness and finite
propagation.  Its local trace is therefore a finite sum.  Since the
operator and coefficient metric are pulled back from the base, the
local transfer commutes with the differential operator.  The traced
wave is thus the downstairs wave with traced initial data, by
uniqueness.  Integrating in time proves that trace commutes with $K$.
On the compact base, a spectral function with value one at zero is the
identity on harmonic representatives and therefore on Dolbeault
cohomology, also computed by currents.  This proves~\textup{(iii)}.

It remains to check the dependence on $S$.  For an $E$-valued $(0,n)$
block of $y$ antiholomorphic degree $q$, the unscaled metric on $E$ gives
\begin{equation}\label{sp:scaling-block}
 |B|_S=S^{-q/2}|B|_1,\qquad
 \dd V_S=S^n\dd V_1,\qquad
 \|B\|_S=S^{(n-q)/2}\|B\|_1.
\end{equation}
Put $a_C^S(B)=\int_C|B|_S\dd V_1$.
The kernel bound and fixed propagation imply
$\|KB\|_1\le C\|(a_C^1(B))_C\|_{\ell^2}$: each output cell sees
only boundedly many input neighbors, and the neighbor matrix has
uniformly bounded row and column sums.  Since $K$ preserves $q$,
\[
 \|KB\|_S
 \le C S^{(n-q)/2}S^{q/2}
             \|(a_C^S(B))_C\|_{\ell^2}
 = C S^{n/2}\|(a_C^S(B))_C\|_{\ell^2}.
\]
The finitely many orthogonal blocks give
\eqref{sp:smoothing-norm}.  Product cells may overlap and have bounded
finite chart multiplicities; these change only the fixed constant.
\end{proof}

\subsection{Cell weights and one localization step}
\label{sp:localization}

We now combine the weighted inverse with the fixed-propagation
smoother. One step replaces a compactly supported closed form by a
smaller representative of the same traced class. The thresholds are
nested so that the output support of one step is an admissible input
for the next.

Fix an integer $r>30n+6$ and choose $0<\zeta<1/(8r)$.  These choices
are made before $N$ tends to infinity.  At step $j\in\{1,\ldots,r\}$ set
\begin{equation}\label{sp:thresholds}
 S=N^2,\quad H_j=N^{3/4-2(j-1)\zeta},\quad
 H_j^-=H_jN^{-\zeta},\quad t_j=\frac{2\log N}{H_j}.
\end{equation}
In particular $H_j\ge N^{1/2}$, $H_j\le S$, and $t_jS\to\infty$.
Constants in the following construction may depend on the fixed $r$.

For each step choose a bump $w^{(j)}\in C_c^\infty(U,[0,1])$, equal
to one on a neighborhood of $L$, and let
$w_a^{(j)}(y)=w^{(j)}(a^{-1}y)$.  The bumps can be chosen as squares so
that
\[
 |\dd w_a^{(j)}|^2\le C_jw_a^{(j)}.
\]
Their translates have bounded overlap and bounded derivatives.  Choose
the support at step $j+1$ sufficiently large that $w^{(j+1)}$ equals
one on every fixed-radius neighborhood of $\supp w^{(j)}$ required
below.  Only finitely many enlargements are made, all independent of
$N$.

We suppress $j$ temporarily and write $H,H^-,t,w_a$ for the data at
that step.  Let $\rho$ be a smooth convex regularization of
$\max(0,s)$, equal to zero for $s\le-1$ and to $s$ for $s\ge1$,
with $0\le\rho'\le1$, $\rho''\ge0$, and $\rho\ge\max(0,s)$.
Define
\begin{align*}
 \ell_a(x)&=-2H+\rho\bigl(u_N(a^{-1}x)+2H\bigr),\\
 p_H(\ell)&=\ell+\frac{\ell^2}{8H},\\
 T(x,y)&=\sum_{a\in P}w_a(y)p_H(\ell_a(x)),\\
 \varphi(x,y)&=-H/2+\rho\bigl(T(x,y)+H/2\bigr).
\end{align*}
For large $N$, the clipped functions satisfy
$-2H\le\ell_a\le0$ and $-H/2\le\varphi\le0$.
The sum is locally finite and invariant under the diagonal action.
Moreover
\begin{equation}\label{sp:weight-curvature}
 \dd\dc\varphi\ge-C_j(\omega_x+H\omega_y).
\end{equation}
Here is the estimate including mixed terms.  On $[-2H,0]$,
$1/2\le p_H'\le1$ and $p_H''=1/(4H)$.  The $x$ Hessian of each
summand contains
$w_ap_H''\ii\partial\ell_a\wedge\dbar\ell_a$ as a positive
term.  The mixed derivative is bounded by half this term plus a
constant times
$H|\dd w_a|^2/w_a$ in the $y$ direction, using Cauchy--Schwarz;
the expression is interpreted as zero where $w_a=0$.  The bump
inequality bounds the latter term by $C_jH\omega_y$.  The pure $y$
Hessian has the same bound because $|p_H(\ell_a)|\le2H$.
The remaining $x$ Hessian is bounded below by $-w_a\omega_x$.
Bounded overlap and the final convex clipping give
\eqref{sp:weight-curvature}.

Put
\begin{equation}\label{sp:actual-weights}
 h_j=e^{-t_j\varphi},\qquad k_j=h_j^{-1},\qquad b=\log N.
\end{equation}
These are smooth multipliers on $D$.  They satisfy
\eqref{sp:weight-hypotheses}, with a constant independent of $N$ and
$S$.  Indeed $1\le h_j\le N$, and
\[
 \ii\Theta_{E^*,h_j}
 =\ii\Theta_{E^*}+t_j\dd\dc\varphi.
\]
The first term is bounded below by $-C\omega_y$, which is absorbed
by $-C't_j\Omega_S$ because $t_jS\to\infty$; the second is controlled
by \eqref{sp:weight-curvature} and $H_j\le S$.

Call the following set the deep region at step $j$:
\begin{equation}\label{sp:deep-region}
 \mathcal D_j=\{(x,y): x\in gL,\ A_a(g)\ge H_j,
                       \ w_a^{(j)}(y)=1
                       \text{ for some }a,g\}/P.
\end{equation}
On this region $\ell_a\le-H_j$ for a witnessing index and
$p_{H_j}(\ell_a)\le-7H_j/8$.  The final clipping is consequently
constant, $\varphi=-H_j/2$, and
\begin{equation}\label{sp:deep-weight}
 h_j=N,\qquad k_j=N^{-1}\quad\text{on }\mathcal D_j.
\end{equation}

A product cell $gL\times g'L$, modulo the diagonal action, is marked
at step $j$ when some $a$ satisfies
\begin{equation}\label{sp:marked-cells}
 A_a(g)\ge H_j^-,\qquad
 g'L\cap\supp w_a^{(j)}\ne\varnothing.
\end{equation}
There are at most $C_jN^{2n}$ marked cells.  In fact, translating by
$a^{-1}$ leaves only a bounded number of possible second cells meeting
$\supp w^{(j)}$, and \eqref{sp:cell-count} bounds the number of first
cells.  Any fixed enlargement of their union has the same bound on its
number of cells and on its standard volume.

On every unmarked cell $C$ one has
\begin{equation}\label{sp:unmarked-weight}
 \int_C h_j\,\dd V_1\le C_j.
\end{equation}
Indeed $\ell_a\ge u_N(a^{-1}x)$, $p_H(\ell_a)\ge\ell_a$, and
$\varphi\ge T$.  Since all the potentials are nonpositive and
$0\le w_a\le1$, this bounds $h_j$ on a cell by the product of
$\exp(-t_ju_N(a^{-1}x))$ over the bounded number of indices whose
$y$ supports meet that cell.  Each such index has $A_a(g)<H_j^-$
when the cell is unmarked.  H\"older's inequality and
\eqref{sp:cell-exponential} apply, because
\[
 t_j(H_j^-+1)\le 2(\log N)(N^{-\zeta}+N^{-1/2})\longrightarrow0.
\]
Integration in the $y$ cell contributes only a fixed volume factor.

Choose a smooth compactly supported cutoff $\chi_j$ equal to one on
a neighborhood of all marked cells, with support in a fixed enlargement
of their union and with $|\dd\chi_j|_1\le C_j$.  Such cutoffs are
obtained by smoothing the distance cutoff in the uniform charts.
Every cell meeting $\supp\dd\chi_j$ is unmarked.
Enlarge the supports of $w^{(j+1)}$ as specified above so that
\begin{equation}\label{sp:nested-supports}
 \text{the $R$-neighborhood of $\supp\chi_j$ lies in }
 \mathcal D_{j+1}\qquad(j<r).
\end{equation}
To verify this, a point in that neighborhood has an $x$ cell within a
fixed distance of a cell witnessing \eqref{sp:marked-cells}.
By \eqref{sp:cell-neighbors}, its depth is at least
$c_jH_j^--1$, which exceeds $H_{j+1}$ for large $N$, since
$H_j^-/H_{j+1}=N^\zeta$.  Its $y$ coordinate lies within a fixed
distance of $\supp w_a^{(j)}$, where the next bump is one.  This
proves \eqref{sp:nested-supports}.  Similarly,
\eqref{sp:potential-depth} shows that $\mathcal D_1$ contains every
prescribed fixed-radius neighborhood of the lifted diagonal for large
$N$.

\begin{lemma}\label{sp:one-step}
Let $F$ be a smooth compactly supported closed $E$-valued $(0,n)$-form
supported in $\mathcal D_j$.  There is a smooth compactly supported
closed form $F_{\mathrm{new}}$ of the same type whose traced
cohomology class is that of $F$, with
\begin{equation}\label{sp:one-step-norm}
 \|F_{\mathrm{new}}\|_S\le C_jN^{-1/6}\|F\|_S.
\end{equation}
Its support is in the fixed $R$-neighborhood of $\supp\chi_j$, has
standard volume at most $C_jN^{2n}$, and lies in $\mathcal D_{j+1}$
when $j<r$.
\end{lemma}

\begin{proof}
By Lemma~\ref{sp:weighted-inverse} and
\eqref{sp:deep-weight}, solve $\dbar v=F$ with
\begin{align*}
 \|v\|_{S,k_j}
 &\le C_j\bigl(N^{1/4}+\sqrt{t_j}\,N^{1/2}\bigr)
                              N^{-1/2}\|F\|_S\\
 &\le C_jN^{-1/6}\|F\|_S.
\end{align*}
The last inequality uses
$\sqrt{t_j}\le C(\log N)^{1/2}N^{-1/4}$.
The compactly supported closed current
\[
 B=F-\dbar(\chi_jv)=-\dbar\chi_j\wedge v
\]
has the same traced cohomology class as $F$.  The equality uses
$\chi_j=1$ near $\supp F$.  On any cell meeting $\supp B$, which is
unmarked, Cauchy--Schwarz and \eqref{sp:unmarked-weight} give
\[
 \int_C|B|_S\,\dd V_1
 \le C_j\left(\int_C|v|_S^2 k_j\,\dd V_1\right)^{1/2}.
\]
Here wedging with $\dbar\chi_j$ has bounded operator norm for
$\Omega_S$, because $S\ge1$.  After summing squares and using bounded
overlap and $\dd V_S=S^n\dd V_1$, we obtain
\[
 \left(\sum_C\left(\int_C|B|_S\,\dd V_1\right)^2\right)^{1/2}
 \le C_jS^{-n/2}N^{-1/6}\|F\|_S.
\]
Set $F_{\mathrm{new}}=KB$ with the fixed operator of
Lemma~\ref{sp:smoothing-lemma}.  Its norm estimate is
\eqref{sp:one-step-norm}; its closedness and traced class follow from
the same lemma.  Fixed propagation, the marked-cell count and
\eqref{sp:nested-supports} give all support assertions.  Only the
compactly supported current $\chi_jv$ is traced; no sum of the global
solution $v$ is used.
\end{proof}

\subsection{The compact diagonal class}
\label{sp:diagonal-class}

\begin{proof}[Proof of Theorem~\ref{thm:zero-spectrum}]
Continue under the assumption of positive lower bounds in all $(p,0)$
degrees.  Let $\iota:M\to D$ be the lifted diagonal and let
$[\iota(M)]$ be its integration current of type $(n,n)$.  Orbifold
integration defines this current directly by pullback of test forms to
$M$, regardless of singularities of the coarse spaces.
Project it to the component whose holomorphic degree lies entirely in
the $y$ factor.  Under the natural identification with $E$-valued
$(0,n)$ currents, call the result $T_0$.  This projection commutes with
$\dbar$, so $T_0$ is closed and compactly supported.

Its trace to $M\times M$ has a nonzero cohomology class.  More
explicitly, the canonical isomorphism
$K_{M\times M}\otimes E_0^*\simeq p_x^*K_M$ identifies
$p_x^*\omega^n$ with a closed $E_0^*$-valued $(m,n)$ test form $\beta$.
The holomorphic-degree projection used to define $T_0$ does not change
the pairing with this form, since all other components vanish on
wedging with $p_x^*\omega^n$.  Hence
\begin{equation}\label{sp:diagonal-pairing}
 \langle\operatorname{Tr}T_0,\beta\rangle
 =\int_M\omega^n>0.
\end{equation}

Let $F_0=KT_0$.  It is smooth, compactly supported, and closed; its
traced class has the same pairing \eqref{sp:diagonal-pairing}.  It is
independent of $N$, and its support lies in a fixed neighborhood of the
lifted diagonal.  Thus it lies in $\mathcal D_1$ for all sufficiently
large $N$.  Apply Lemma~\ref{sp:one-step} successively for the fixed
number $r$ of steps to obtain $F_r$.  Its standard support volume is
at most $C_rN^{2n}$, while
\[
 \|F_r\|_S\le C_rN^{-r/6}\|F_0\|_S.
\]
For degree $(0,n)$, \eqref{sp:scaling-block} and $0\le q\le n$ show
that $\|F_r\|_1\le\|F_r\|_S$ and
$\|F_0\|_S\le S^{n/2}\|F_0\|_1$.  The constant norm conversions
are used only at these endpoints.  Since $S=N^2$,
\[
 \|F_r\|_1\le C_rN^{n-r/6}.
\]
The pullback of the fixed test form $\beta$ is uniformly bounded for
the standard metric.  Therefore Cauchy--Schwarz gives
\[
 \bigl|\langle\operatorname{Tr}F_r,\beta\rangle\bigr|
 \le C\vol_1(\supp F_r)^{1/2}\|F_r\|_1
 \le C_rN^{2n-r/6}\longrightarrow0.
\]
Our fixed choice $r>30n+6$ more than suffices for this decay.  Every
step preserves the traced cohomology class, so its pairing is the
strictly positive constant in \eqref{sp:diagonal-pairing}.  This is a
contradiction.  At least one scalar $(p,0)$ Dolbeault Laplacian has no
positive lower bound, which proves the theorem.
\end{proof}
